\chapter{Lecture 12: The Schwarzschild Horizon and Black-Hole Interior}

\lectureoverview{The geometry near a Schwarzschild horizon is first
recognized in a simpler setting: flat spacetime described by uniformly
accelerated observers.  Hyperbolic coordinates produce the Rindler
metric, an observer-dependent horizon, and the asymmetry between an
inertial observer who crosses that horizon and an accelerated observer
who remains outside.  The Schwarzschild metric appears to be singular
at \(r=2GM\), but a proper-distance coordinate and a rescaled time
reduce its near-horizon form exactly to the Rindler metric.  The
horizon is therefore not a curvature singularity.  It is nevertheless
a causal boundary for observers who stay outside.  The true
Schwarzschild singularity is at \(r=0\), where tidal curvature
diverges and which lies as a spacelike future boundary inside the
hole.  The closing discussion returns to the vacuum Einstein
equations, the Newtonian exterior problem, gravitational collapse,
mean-density scaling, cosmic horizons, and the limits of classical
physics at the singularity.}

We use signature \(+---\), set \(c=1\), and write
\begin{equation}
  r_{\mathrm s}=2GM
  \label{eq:l12-rs}
\end{equation}
for the Schwarzschild radius.  The symbol \(\rho\) will denote proper
distance from a horizon; \(r\) is reserved for the Schwarzschild
areal radius.

\section{The Accelerated Congruence}

Before entering a black hole, return to flat spacetime.  In an
ordinary inertial diagram, every observer in the accelerated
congruence follows a hyperbola.  The curves nearer the null boundary
turn through a large rapidity in a shorter proper distance and
therefore require greater proper acceleration.  This differs from the
Newtonian picture in which a rigid array is assigned one common
acceleration.  Relativity permits a stationary accelerated array only
when its proper acceleration varies from one worldline to the next.

\begin{classroomqa}
\audiencequestion{Does this show that special relativity cannot be
used for accelerated observers?
\lecturetimestamp{00:01:35--00:02:54}}
\lecturerresponse{No.  The entire construction takes place in flat
Minkowski spacetime and is governed by special relativity.  Each
fixed-position observer is accelerated, and different positions
require different accelerations.  That variation is also what one
expects for observers supported at different heights in a
gravitational field.}
\end{classroomqa}

Introduce inertial coordinates \((T,X,y,z)\) and hyperbolic
coordinates \((\omega,\rho,y,z)\):
\begin{equation}
  X=\rho\cosh\omega,
  \qquad
  T=\rho\sinh\omega.
  \label{eq:l12-rindler-map}
\end{equation}
These equations are the Lorentzian counterparts of
\(x=\rho\cos\theta\), \(y=\rho\sin\theta\).  They imply
\begin{equation}
  X^2-T^2=\rho^2,
  \qquad
  \tanh\omega=\frac{T}{X}.
  \label{eq:l12-rindler-curves}
\end{equation}
Thus \(\rho=\mathrm{constant}\) labels hyperbolic worldlines and
\(\omega=\mathrm{constant}\) labels straight rays through the origin.
The rapidity \(\omega\) ranges from \(-\infty\) to \(+\infty\).

\begin{figure}[tbp]
  \centering
  \includegraphics[width=0.84\textwidth]{lecture_12_figure_01.png}
  \caption{The Rindler congruence, constant-\(\omega\) rays, and the
  coordinate map \(X=\rho\cosh\omega\),
  \(T=\rho\sinh\omega\).}
  \lectureframe{00:09:55}
\end{figure}

\subsection{What Alice and Bob See}

The null lines \(X=\pm T\) divide Minkowski spacetime into wedges.  An
observer confined to the right Rindler wedge cannot receive a signal
emitted beyond its past or future null boundary.  The future boundary
\(X=T\) is approached as \(\omega\to\infty\).  It is an acceleration
horizon, not a black-hole horizon: there is no curvature and no
gravitating body.

Let Alice release herself and cross the null boundary inertially.  Let
Bob remain on a fixed-\(\rho\) hyperbola.  Seeing means receiving
light, so Bob locates Alice by tracing past-directed null rays from
successive events on his worldline.  Each later ray meets Alice closer
to the horizon.  Bob therefore sees her approach the boundary
asymptotically.  Alice crosses it after finite proper time.

\begin{figure}[tbp]
  \centering
  \includegraphics[width=0.80\textwidth]{lecture_12_figure_02.png}
  \caption{Past-directed null rays used by Bob to see Alice.
  Successive reception events intersect Alice's inertial worldline
  closer to the Rindler horizon.}
  \lectureframe{00:12:55}
\end{figure}

Bob sees Alice's processes slow and her longitudinal dimensions
Lorentz contract as their relative rapidity grows.  Alice does not see
a wall or a singular event.  Looking backward, she sees Bob
accelerating away.  His signals are also Doppler shifted and slowed,
but he does not disappear behind a horizon in her causal past.

\begin{classroomqa}
\audiencequestion{Does Bob see Alice's time slow, and what does Alice
see when she looks back at Bob?
\lecturetimestamp{00:11:04--00:12:59}}
\lecturerresponse{Bob receives Alice's successive signals farther
apart and says that her clock slows as she approaches his horizon.
Alice sees Bob receding on an accelerated trajectory.  She can
continue to receive signals from him; she does not assign him the same
one-way disappearance that he assigns to her.  The asymmetry comes
from Bob's acceleration and from the different null rays available to
the two observers.}
\end{classroomqa}

\subsection{Clock Rates in the Accelerated Frame}

Substitution into the Minkowski line element gives
\begin{align}
  ds^2
  &=dT^2-dX^2-dy^2-dz^2 \notag\\
  &=\rho^2d\omega^2-d\rho^2-dy^2-dz^2.
  \label{eq:l12-rindler-metric}
\end{align}
For an observer at fixed \(\rho,y,z\),
\begin{equation}
  d\tau=\rho\,d\omega.
  \label{eq:l12-rindler-proper-time}
\end{equation}
The coordinate \(\omega\) is a common time label, not a common proper
time.  Identical proper-time clocks at different \(\rho\) advance at
different rates per unit \(\omega\).  Conversely, clocks engineered to
display \(\omega\) must be calibrated differently at different
positions.

The proper acceleration of a fixed-\(\rho\) observer is
\begin{equation}
  a(\rho)=\frac{1}{\rho}.
  \label{eq:l12-rindler-acceleration}
\end{equation}
The acceleration diverges as the observer approaches \(\rho=0\), the
Rindler horizon.

\begin{figure}[tbp]
  \centering
  \includegraphics[width=0.82\textwidth]{lecture_12_figure_03.png}
  \caption{The board relation \(d\tau=\rho\,d\omega\), showing why
  equal increments of accelerated time do not represent equal proper
  time at different \(\rho\).}
  \lectureframe{00:14:50}
\end{figure}

\begin{classroomqa}
\audiencequestion{If the transformation is only a change of
coordinates, where does the acceleration enter?
\lecturetimestamp{00:13:02--00:13:57}}
\lecturerresponse{The map itself is only a coordinate transformation.
The physical statement comes when observers are placed on
\(\rho=\mathrm{constant}\) curves.  Those curves are hyperbolae in the
inertial diagram, so the observers occupying them accelerate.  A
coordinate chart does not accelerate anyone; a chosen congruence of
worldlines does.}
\end{classroomqa}

\begin{classroomqa}
\audiencequestion{Is \(d\omega\) the proper-time increment on every
fixed-\(\rho\) worldline?
\lecturetimestamp{00:13:59--00:15:30}}
\lecturerresponse{No.  Equation
\eqref{eq:l12-rindler-proper-time} gives
\(d\tau=\rho\,d\omega\).  The factor of \(\rho\) is constant along
one worldline but differs between worldlines.  A common coordinate
time can be imposed, but it is not the reading of identical local
clocks everywhere.}
\end{classroomqa}

\subsection{Equal acceleration is not a rigid accelerated frame}

A translated copy of one hyperbola has the same proper acceleration
as the original.  It does not, however, belong to the stationary
Rindler array.  Equal coordinate separation in the original inertial
frame is not equal separation on the instantaneous rest slices of the
accelerated observers.  Lorentz contraction changes the latter as the
motion develops.  The fixed-\(\rho\) family instead changes
acceleration with position so that neighboring observers retain a
constant radar distance.

\begin{figure}[tbp]
  \centering
  \includegraphics[width=0.80\textwidth]{lecture_12_figure_04.png}
  \caption{Translated equal-acceleration hyperbolae compared with the
  stationary Rindler congruence.  The inertial separations may match
  while the observers' instantaneous rest-frame separations do not.}
  \lectureframe{00:18:55}
\end{figure}

\begin{classroomqa}
\audiencequestion{Would two translated hyperbolae of the same shape
have different gravitational potentials or forces?
\lecturetimestamp{00:15:51--00:19:50}}
\lecturerresponse{The observers have the same proper acceleration and
feel the same support force.  An additive constant in a Newtonian
potential is not observable.  What fails is rigidity: translated
hyperbolae do not remain at fixed mutual distance when distance is
measured on their changing rest slices.  The Rindler congruence uses
different accelerations precisely to avoid that time-dependent
separation.}
\end{classroomqa}

\section{The Schwarzschild Exterior}

The accelerated-frame picture becomes useful because the
Schwarzschild horizon has the same local geometry.  In Schwarzschild
coordinates,
\begin{equation}
  \begin{aligned}
    ds^2
    &=
    \left(1-\frac{2GM}{r}\right)dt^2
    -
    \frac{dr^2}{1-\frac{2GM}{r}}\\
    &\quad
    -
    r^2d\Omega^2,\\
    d\Omega^2
    &=
    d\theta^2+\sin^2\theta\,d\phi^2.
  \end{aligned}
  \label{eq:l12-schwarzschild}
\end{equation}
This is the exterior metric of every spherically symmetric body.  It
describes the field outside the Earth or a star as well as the vacuum
region outside a black hole.  For an extended body it need not hold
through the matter; after complete collapse the vacuum solution
extends down to the singular region.

\begin{figure}[tbp]
  \centering
  \includegraphics[width=0.86\textwidth]{lecture_12_figure_05.png}
  \caption{The Schwarzschild line element on the board.  The complete
  normalized equation is given in
  \eqref{eq:l12-schwarzschild}.}
  \lectureframe{00:22:30}
\end{figure}

Two radii look suspicious.  At \(r=r_{\mathrm s}\), the coefficient of
\(dt^2\) vanishes while the coefficient of \(dr^2\) diverges.  At
\(r=0\), the coefficients also diverge or vanish and their signs are
reversed relative to the exterior.  A divergent metric component does
not by itself prove a physical singularity.  We must ask whether an
invariant curvature or a freely falling measurement diverges.

The distinction is decisive:
\begin{align}
  R_{\mu\nu}&=0,
  & R&=0
  &&\text{for }r>0, \notag\\
  R_{\alpha\beta\gamma\delta}
  R^{\alpha\beta\gamma\delta}
  &=\frac{48G^2M^2}{r^6}.
  \label{eq:l12-kretschmann}
\end{align}
The Kretschmann scalar is finite at \(r=2GM\) and diverges at
\(r=0\).  The horizon singularity is a coordinate failure; the central
singularity is curvature.

\section{A Coordinate Singularity Manufactured in Flat Space}

The Rindler metric can be made to look as alarming as the Schwarzschild
metric.  Define an auxiliary coordinate
\begin{equation}
  R_{\mathrm{aux}}=\rho^2.
  \label{eq:l12-r-aux}
\end{equation}
Then
\begin{equation}
  d\rho^2=\frac{dR_{\mathrm{aux}}^2}{4R_{\mathrm{aux}}},
\end{equation}
and flat spacetime becomes
\begin{equation}
  ds^2
  =
  R_{\mathrm{aux}}\,d\omega^2
  -
  \frac{dR_{\mathrm{aux}}^2}{4R_{\mathrm{aux}}}
  -
  dy^2-dz^2.
  \label{eq:l12-flat-bad-coordinate}
\end{equation}
At \(R_{\mathrm{aux}}=0\), one coefficient vanishes and another
diverges.  No curvature has appeared.  The singularity was created
entirely by replacing \(\rho\) with \(\rho^2\).

\begin{figure}[tbp]
  \centering
  \includegraphics[width=0.84\textwidth]{lecture_12_figure_06.png}
  \caption{The flat Rindler metric rewritten with
  \(R_{\mathrm{aux}}=\rho^2\).  Its zero and pole at
  \(R_{\mathrm{aux}}=0\) imitate the Schwarzschild coordinate
  singularity without introducing curvature.}
  \lectureframe{00:32:40}
\end{figure}

The apparent exchange of space and time is equally harmless in this
example.  Since
\begin{equation}
  X^2-T^2=R_{\mathrm{aux}},
\end{equation}
positive \(R_{\mathrm{aux}}\) labels right-wedge timelike hyperbolae.
For negative \(R_{\mathrm{aux}}\),
\begin{equation}
  T^2-X^2=-R_{\mathrm{aux}},
\end{equation}
and the same coordinate labels upper-wedge spacelike hyperbolae.
Nothing rotates the physical axes.  The coordinate grid changes
character across its null boundary.

\begin{figure}[tbp]
  \centering
  \includegraphics[width=0.82\textwidth]{lecture_12_figure_07.png}
  \caption{Constant-\(R_{\mathrm{aux}}\) hyperbolae on opposite sides
  of the Rindler horizon.  The sign change turns fixed-coordinate
  curves from timelike to spacelike without changing the underlying
  flat geometry.}
  \lectureframe{00:37:55}
\end{figure}

\begin{classroomqa}
\audiencequestion{How can \(\rho^2\) become negative?
\lecturetimestamp{00:34:25--00:38:30}}
\lecturerresponse{The coordinate \(\rho\) was defined only in the
right wedge, so its square is nonnegative there.  The analytic
continuation is made with \(R_{\mathrm{aux}}=X^2-T^2\).  That
quantity is negative in the future wedge, where the level sets are
spacelike hyperbolae.  It is the continued coordinate, not a real
square of a real right-wedge distance, that changes sign.}
\end{classroomqa}

\section{The Near-Horizon Approximation}

Write
\begin{equation}
  1-\frac{r_{\mathrm s}}{r}
  =\frac{r-r_{\mathrm s}}{r}.
  \label{eq:l12-factor}
\end{equation}
We examine a small neighborhood of one point on the horizon.  The
difference \(r-r_{\mathrm s}\) must be retained because it is small,
vanishes at the horizon, and changes sign across it.  Factors of \(r\)
that vary only by a tiny fractional amount may be replaced by
\(r_{\mathrm s}\).  Therefore
\begin{equation}
  ds^2
  \simeq
  \frac{r-r_{\mathrm s}}{r_{\mathrm s}}\,dt^2
  -
  \frac{r_{\mathrm s}}{r-r_{\mathrm s}}\,dr^2
  -
  r_{\mathrm s}^2d\Omega^2.
  \label{eq:l12-near-horizon-before-plane}
\end{equation}

The angular part is a sphere of radius \(r_{\mathrm s}\).  Under a
microscope, a small patch of that sphere is its tangent plane.  With
local Cartesian coordinates \(y,z\),
\begin{equation}
  r_{\mathrm s}^2d\Omega^2\simeq dy^2+dz^2,
\end{equation}
and hence
\begin{equation}
  ds^2
  \simeq
  \frac{r-r_{\mathrm s}}{r_{\mathrm s}}\,dt^2
  -
  \frac{r_{\mathrm s}}{r-r_{\mathrm s}}\,dr^2
  -
  dy^2-dz^2.
  \label{eq:l12-near-horizon-local}
\end{equation}

\begin{figure}[tbp]
  \centering
  \includegraphics[width=0.82\textwidth]{lecture_12_figure_08.png}
  \caption{A small patch of the horizon two-sphere replaced by its
  tangent plane, so
  \(r_{\mathrm s}^2d\Omega^2\simeq dy^2+dz^2\).}
  \lectureframe{00:45:10}
\end{figure}

\begin{classroomqa}
\audiencequestion{Is the useful approximation determined by how far
one is from the horizon?
\lecturetimestamp{00:38:36--00:41:55}}
\lecturerresponse{Yes, and also by how small a patch of the horizon is
examined.  Near \(r=r_{\mathrm s}\), slowly varying factors may be
evaluated at \(r_{\mathrm s}\), but the difference
\(r-r_{\mathrm s}\) must remain.  On a sufficiently small angular
patch the horizon sphere is also well approximated by a plane.}
\end{classroomqa}

\subsection{Proper distance from the horizon}

At fixed \(t,y,z\), define \(\rho\) as radial proper distance from the
horizon:
\begin{equation}
  d\rho
  =
  \sqrt{\frac{r_{\mathrm s}}{r-r_{\mathrm s}}}\,dr.
  \label{eq:l12-drho}
\end{equation}
Integrating from \(r_{\mathrm s}\) to an exterior point \(r\),
\begin{align}
  \rho(r)
  &=
  \int_{r_{\mathrm s}}^r
  \sqrt{\frac{r_{\mathrm s}}{r'-r_{\mathrm s}}}\,dr'
  \notag\\
  &=
  2\sqrt{r_{\mathrm s}(r-r_{\mathrm s})}
  =
  \sqrt{8GM(r-2GM)}.
  \label{eq:l12-rho-result}
\end{align}
Equivalently,
\begin{equation}
  r-r_{\mathrm s}
  =
  \frac{\rho^2}{4r_{\mathrm s}}
  =
  \frac{\rho^2}{8GM}.
  \label{eq:l12-r-minus-rs}
\end{equation}

\begin{figure}[tbp]
  \centering
  \includegraphics[width=0.78\textwidth]{lecture_12_figure_09.png}
  \caption{The completed proper-distance relation
  \(\rho=\sqrt{8GM(r-2GM)}\).}
  \lectureframe{00:56:25}
\end{figure}

\begin{classroomqa}
\audiencequestion{Why are the other coordinate increments set to zero,
and what is the upper limit of the proper-distance integral?
\lecturetimestamp{00:48:44--00:53:57}}
\lecturerresponse{We want the distance measured along a radial curve
on a constant-time slice, so \(dt=dy=dz=0\).  The lower endpoint is
the horizon \(r'=r_{\mathrm s}\); the upper endpoint is the exterior
point whose Schwarzschild coordinate is \(r\).  A primed \(r'\) is
used only as the integration variable.  Taking the square root of the
radial metric coefficient supplies the factor
\(\sqrt{r_{\mathrm s}}\) that must remain in the result.}
\end{classroomqa}

Use \eqref{eq:l12-r-minus-rs} in
\eqref{eq:l12-near-horizon-local}.  The radial term is
\(-d\rho^2\) by construction, and the time term becomes
\begin{equation}
  \frac{r-r_{\mathrm s}}{r_{\mathrm s}}dt^2
  =
  \frac{\rho^2}{4r_{\mathrm s}^2}dt^2.
\end{equation}
Now define the dimensionless accelerated time
\begin{equation}
  \omega=\frac{t}{2r_{\mathrm s}}
  =\frac{t}{4GM},
  \qquad
  d\omega=\frac{dt}{4GM}.
  \label{eq:l12-omega}
\end{equation}
The metric becomes
\begin{equation}
  \boxed{
  ds^2
  \simeq
  \rho^2d\omega^2-d\rho^2-dy^2-dz^2.}
  \label{eq:l12-rindler-result}
\end{equation}
Thus the near-horizon Schwarzschild geometry has exactly the Rindler
form given in \eqref{eq:l12-rindler-metric}.  Dividing Schwarzschild
time by \(4GM\) turns it into the corresponding hyperbolic angle.

\begin{figure}[tbp]
  \centering
  \includegraphics[width=0.82\textwidth]{lecture_12_figure_10.png}
  \caption{The final rescaling
  \(\omega=t/(4GM)\) and the near-horizon Rindler metric.}
  \lectureframe{01:00:50}
\end{figure}

\begin{classroomqa}
\audiencequestion{Why may the radial term simply be replaced by
\(d\rho^2\), rather than replacing \(dr\) by \(d\rho\)?
\lecturetimestamp{01:03:38--01:05:35}}
\lecturerresponse{The definition is
\(d\rho^2=[r_{\mathrm s}/(r-r_{\mathrm s})]dr^2\).
It is the complete weighted radial term that equals \(d\rho^2\);
one never asserts \(d\rho=dr\).  The coordinate was chosen precisely
to absorb that metric coefficient.}
\end{classroomqa}

The coefficient
\begin{equation}
  \kappa=\frac{1}{4GM}
  \label{eq:l12-surface-gravity}
\end{equation}
is the Schwarzschild surface gravity in units \(c=1\).  In the local
form \(ds^2=(\kappa\rho)^2dt^2-d\rho^2-\cdots\), a supported observer
at proper distance \(\rho\) has leading acceleration \(a\simeq1/\rho\).
Far from the horizon, Schwarzschild \(t\) again behaves as ordinary
inertial time.  Close to the horizon, \(t/(4GM)\) behaves as Rindler
rapidity.  The transition belongs to the coordinates; the local
geometry at the horizon is regular.

\section{The Horizon as Seen from Outside}

The near-horizon reduction now tells us how an observer who remains
outside describes infall.  A stationary Schwarzschild observer is the
black-hole analogue of a Rindler observer: staying at fixed \(r\)
requires support, and the required acceleration grows without bound as
\(r\to r_{\mathrm s}^{+}\).  For that exterior system of clocks, the
horizon lies at
\[
  \omega=\frac{t}{4GM}\longrightarrow+\infty .
\]
Nothing reaches the horizon at a finite value of Schwarzschild time.
This is not a statement about the infalling observer's proper time.
It is a statement about the time coordinate tied to observers who
refuse to fall.

\begin{classroomqa}
\audiencequestion{How does the infinite Rindler time translate into
the black-hole picture?
\lecturetimestamp{01:10:13--01:12:59}}
\lecturerresponse{The horizon is the
\(\omega=\infty\) boundary of the accelerated chart.  A stationary
observer can wait arbitrarily long, look back, and still assign the
infalling body a position outside the horizon.  Very close to the
horizon this is exactly the Rindler description.  Far away the metric
returns to ordinary asymptotically flat coordinates, so the
hyperbolic interpretation of \(t/(4GM)\) is only local to the
horizon.}
\end{classroomqa}

How, then, can a black hole grow if the exterior observer never sees
matter cross?  In the exterior description, the infalling energy
becomes concentrated ever closer to the old horizon.  Its
gravitational field changes the geometry, and the new horizon encloses
it.  A fully dynamical account requires solving Einstein's equations
with infalling stress-energy; the fixed Schwarzschild metric alone
cannot describe the growth.

\begin{classroomqa}
\audiencequestion{If nothing is seen to cross, how can the black hole
become larger?
\lecturetimestamp{01:13:16--01:14:39}}
\lecturerresponse{The incoming energy is deposited arbitrarily close
to the horizon and changes the spacetime.  The horizon grows outward
and merges with that energy.  Treating the matter as a test object in
an eternally fixed metric misses precisely this backreaction.}
\end{classroomqa}

The optical signature follows directly from the null rays.  Suppose
Alice carries an atomic dipole that emits one oscillation after
another at a fixed frequency in her local frame.  Successive crests
reach Bob at increasingly separated values of his proper time.  In
the local Rindler picture this is an ordinary Doppler shift caused by
their increasing relative rapidity; in Schwarzschild language it is
the gravitational redshift.  These are two descriptions of the same
near-horizon effect.

\begin{figure}[tbp]
  \centering
  \includegraphics[width=0.84\textwidth]{lecture_12_figure_11.png}
  \caption{Successive wave crests from the infalling Alice arrive
  farther apart on Bob's supported worldline.  The same construction
  is read as a Doppler shift in Rindler coordinates and as
  gravitational redshift in Schwarzschild coordinates.}
  \lectureframe{01:17:00}
\end{figure}

Bob first receives optical light, then infrared, microwaves, and radio
waves of progressively longer wavelength.  He receives photons less
frequently, and each photon has lower energy,
\begin{equation}
  E_{\mathrm{obs}}=\hbar\omega_{\mathrm{obs}},
  \qquad
  \omega_{\mathrm{obs}}\longrightarrow0 .
  \label{eq:l12-redshift-energy}
\end{equation}
Every visible process on Alice appears to slow for the same reason:
all information reaches Bob through these increasingly delayed
signals.

\begin{classroomqa}
\audiencequestion{Why begin with optical light rather than radio
signals?
\lecturetimestamp{01:14:53--01:15:24}}
\lecturerresponse{Beginning at a short wavelength makes the redshift
sequence visible in familiar bands: optical, infrared, microwave, and
then successively longer radio waves.  Starting with radio waves would
leave no equally familiar names for the increasingly long
wavelengths.}
\end{classroomqa}

For an emitter with finite proper frequency, only finitely many
oscillations fit into any finite interval of Alice's proper time.
There is therefore a last crest emitted before crossing within that
idealized pulse train, and Bob receives it only after an enormous
delay and redshift.  Increasing the source frequency inserts more
crests, but does not change the limiting geometry.

\begin{classroomqa}
\audiencequestion{Is the longest possible wavelength necessarily of
order the black-hole radius?
\lecturetimestamp{01:18:58--01:20:29; 02:16:54--02:17:11}}
\lecturerresponse{That conclusion does not follow from the classical
redshift diagram alone.  A finite-frequency source has a last crest
before crossing, but classical theory permits arbitrarily low
observed frequencies.  A physical cutoff or a statement about one
photon per mode requires quantum mechanics, so the proposed
maximum-wavelength argument is set aside.  Alice, meanwhile, sees Bob
recede toward the speed of light and sees his processes Doppler
slowed.}
\end{classroomqa}

\section{Communication Across the Horizon}

The horizon is not locally violent.  Put Alice and Shirley into free
fall side by side.  They remain in ordinary local communication as
they cross, provided their separation is small compared with the
curvature scale.  Neither sees the other freeze.  The freezing belongs
to Bob's supported exterior description, not to the freely falling
laboratory.

\begin{classroomqa}
\audiencequestion{Is the mass really collapsed to a point in the
exterior solution?
\lecturetimestamp{01:21:09--01:22:13}}
\lecturerresponse{The Schwarzschild metric specifies the vacuum
geometry outside a spherical source; it does not by itself say that
ordinary matter has always occupied \(r=0\).  For a star, the
Schwarzschild exterior ends at the stellar surface.  After black-hole
formation, the same vacuum form covers the exterior while the
infalling matter and the singular interior require the global
collapse geometry.}
\end{classroomqa}

\begin{figure}[tbp]
  \centering
  \includegraphics[width=0.82\textwidth]{lecture_12_figure_12.png}
  \caption{Two freely falling companions cross the horizon together.
  Their local exchange of signals remains regular even though a
  supported exterior observer receives Alice's signals ever more
  slowly.}
  \lectureframe{01:24:50}
\end{figure}

\begin{classroomqa}
\audiencequestion{If Alice and Shirley fall together, does either one
see the other disappear or become ionized at the horizon?
\lecturetimestamp{01:22:36--01:25:43}}
\lecturerresponse{No local event marks the crossing.  Nearby freely
falling observers continue to exchange light normally.  A supported
observer attributes extreme redshift and delay to Alice; Alice does
not encounter a hot membrane in this classical discussion.  Any
claim about ionization involves the observer, the state of the
radiation, and quantum physics, not a classical curvature singularity
at the horizon.}
\end{classroomqa}

Nor can Alice evade the interior merely by accelerating sideways.
Near the horizon the transverse directions remain ordinary:
\begin{equation}
  ds^2\simeq(\kappa\rho)^2dt^2-d\rho^2-dy^2-dz^2 .
  \label{eq:l12-local-transverse}
\end{equation}
At fixed \(\rho\), a finite transverse velocity must still satisfy
\((dy/dt)^2+(dz/dt)^2<(\kappa\rho)^2\).  The allowed coordinate speed
shrinks as \(\rho\to0\).  A worldline that tries to skim the horizon
at fixed \(r\) would need unbounded acceleration; once inside, fixed
\(r\) is not a timelike option at all.

\begin{classroomqa}
\audiencequestion{Could an infalling observer accelerate tangentially
and travel along the horizon instead of toward the singularity?
\lecturetimestamp{01:27:09--01:31:54}}
\lecturerresponse{A massive observer must remain inside the local
light cone.  Exactly on the horizon the outward null direction is the
horizon itself; a timelike observer cannot remain there.  Inside, all
future-directed timelike and null directions lead toward smaller
\(r\).  Sideways motion changes where the observer arrives, not
whether the future singularity is reached.}
\end{classroomqa}

\section{Inside the Horizon: Fixed Radius Becomes Spacelike}

The apparent singularity at \(r=r_{\mathrm s}\) disappears in a
regular coordinate system, but the singularity at \(r=0\) does not.
The Kretschmann scalar in \eqref{eq:l12-kretschmann} is finite at
\(r=2GM\) and diverges at \(r=0\).

The reversal of coordinate roles inside the horizon is equally
important.  In the Schwarzschild metric, the factor
\(1-r_{\mathrm s}/r\) is negative for \(r<r_{\mathrm s}\).
Consequently the \(r\)-direction contributes with the timelike sign
while the \(t\)-direction contributes with a spacelike sign.  A line
of fixed \(r\), which represented a supported observer outside, is
spacelike inside.  It cannot be the worldline of matter.

\begin{figure}[tbp]
  \centering
  \includegraphics[width=0.82\textwidth]{lecture_12_figure_13.png}
  \caption{The maximally extended local picture: the horizon is a
  regular null boundary, whereas \(r=0\) is a spacelike future
  singularity reached by every future-directed interior worldline.}
  \lectureframe{01:33:25}
\end{figure}

Calling \(r=0\) spacelike means that it is not a place waiting at the
center in the ordinary Newtonian sense.  Inside the black hole it is
more like a future time: just as an exterior observer cannot avoid
tomorrow, an interior observer cannot remain at fixed \(r\) and avoid
the singularity.  Every future-directed timelike path reaches it in a
finite proper time.

\begin{classroomqa}
\audiencequestion{Does a freely falling observer experience any
special event at the horizon, and how long remains before the
singularity?
\lecturetimestamp{01:32:17--01:34:45}}
\lecturerresponse{Nothing locally singular occurs at crossing; the
strange behavior belongs to Schwarzschild coordinates.  The proper
time from the horizon to \(r=0\) is finite and is of order the
black-hole radius in units \(c=1\), hence of order \(GM\).  The exact
coefficient depends on the infalling geodesic.}
\end{classroomqa}

The radial tidal scale is set by
\begin{equation}
  \left|R_{\hat 0\hat r\hat 0\hat r}\right|
  \sim\frac{GM}{r^3}.
  \label{eq:l12-tidal-scale}
\end{equation}
At the horizon this scales as
\[
  \frac{GM}{r_{\mathrm s}^3}
  =\frac{1}{8G^2M^2}.
\]
Thus a stellar-mass black hole can produce destructive horizon-scale
tidal forces, whereas the horizon of a supermassive black hole can be
locally gentle.  In either case the tidal field diverges as
\(r\to0\).

\begin{classroomqa}
\audiencequestion{Is the singularity an impact with matter, and does
a larger black hole make the crossing safer?
\lecturetimestamp{01:35:23--01:40:21}}
\lecturerresponse{The singularity is not an ordinary material surface
that one strikes.  It is the boundary where the classical curvature
diverges and the theory ceases to predict.  Increasing \(M\) reduces
the curvature and tidal force at the horizon as \(M^{-2}\), so a very
large black hole may be crossed without immediate discomfort.  The
traveler still reaches the high-curvature region later in finite
proper time.}
\end{classroomqa}

\section{Rods, Clocks, and Coordinate Roles}

Coordinates are labels, not instruments.  Outside the horizon, a
static lattice of rods and supported clocks can give operational
meaning to Schwarzschild \(r\) and \(t\).  That lattice cannot be
continued through the horizon: the required support acceleration
diverges there, and fixed \(r\) becomes spacelike inside.  One may
continue the coordinate symbols algebraically, but one must not
continue their exterior measuring instructions unchanged.

\begin{classroomqa}
\audiencequestion{If \(r\) becomes timelike and \(t\) spacelike, do
rods become clocks and clocks become rods?
\lecturetimestamp{01:40:30--01:45:25}}
\lecturerresponse{No physical object changes its nature.  The signs
say that the coordinate directions have exchanged causal roles, not
that a meter stick has turned into a clock.  Inside, \(r\) cannot be
measured by a chain of stationary rulers because no stationary
observers exist there.  A freely falling local laboratory keeps its
ordinary rods, clocks, and Minkowski metric while it crosses.}
\end{classroomqa}

\section{Where the Exterior Metric Comes From}

The Schwarzschild metric is not obtained by placing a point mass into
the line element by hand.  Begin with Einstein's equation,
\begin{equation}
  G_{\mu\nu}
  \equiv
  R_{\mu\nu}-\frac{1}{2}g_{\mu\nu}R
  =8\pi G T_{\mu\nu}.
  \label{eq:l12-einstein}
\end{equation}
Outside a spherical distribution of matter, \(T_{\mu\nu}=0\), so
\begin{equation}
  G_{\mu\nu}=0.
  \label{eq:l12-vacuum}
\end{equation}
Spherical symmetry and asymptotic flatness then select the
Schwarzschild family.  Birkhoff's theorem adds the striking fact that
the spherically symmetric vacuum exterior is static even when the
source itself changes spherically with time.  The integration
constant is the total mass \(M\) enclosed by the vacuum region.

\begin{figure}[tbp]
  \centering
  \includegraphics[width=0.84\textwidth]{lecture_12_figure_14.png}
  \caption{The vacuum Einstein equation and its Newtonian analogue.
  The exterior field is source-free locally, while its integration
  constant records the mass enclosed by the vacuum region.}
  \lectureframe{01:47:55}
\end{figure}

The Newtonian analogue makes the logic transparent:
\begin{equation}
  \nabla^2\Phi=4\pi G\rho_{\mathrm m},
  \qquad
  \nabla^2\Phi=0
  \quad\hbox{outside the matter}.
  \label{eq:l12-poisson}
\end{equation}
For a spherically symmetric potential that vanishes at infinity, the
exterior solution is
\begin{equation}
  \Phi(r)=-\frac{GM}{r}.
  \label{eq:l12-newton-potential}
\end{equation}
The differential equation is homogeneous in the exterior, yet the
boundary condition carries the source mass.  General relativity
works in the same structural way, although its field equations are
nonlinear and the unknown is the metric rather than a scalar
potential.

\begin{classroomqa}
\audiencequestion{What matter distribution was assumed when the
Schwarzschild metric was derived?
\lecturetimestamp{01:46:19--01:50:24}}
\lecturerresponse{No detailed interior distribution is needed for the
vacuum exterior.  One assumes spherical symmetry, solves
\(G_{\mu\nu}=0\) outside a sphere containing all the matter, and
requires asymptotic flatness.  The enclosed mass fixes the one
remaining constant, just as it fixes the coefficient
\(-GM\) in the Newtonian solution \(\Phi=-GM/r\).}
\end{classroomqa}

Shrinking the enclosing sphere in Newtonian theory leaves a
source-free solution everywhere except at the point source.  In a
collapsed Schwarzschild geometry, however, saying that the matter is
at \(r=0\) must be interpreted causally: \(r=0\) is the spacelike
future singularity, not a persistent point in space on which matter
sits while the exterior evolves.

\begin{classroomqa}
\audiencequestion{Where is the material after it falls into the black
hole?
\lecturetimestamp{01:50:31--01:52:26}}
\lecturerresponse{If one insists on locating it in the ideal
Schwarzschild solution, its future ends at \(r=0\).  Everywhere away
from the collapsing matter and the singular boundary, the vacuum
equation holds.  The analogy with a Newtonian point source is useful
for the exterior field, but it fails geometrically because the
relativistic \(r=0\) is spacelike rather than a timelike point at the
center.}
\end{classroomqa}

\section{Why Collapse Is Not Newtonian Point Aiming}

In Newtonian many-body mechanics, collapse to an exact point is
exceptional.  The slightest angular momentum generally makes
particles miss one another and pass through or orbit.  Even if a
point singularity were postulated, an additional particle would hit
it only under exquisitely aimed initial conditions.  This intuition
once encouraged the belief that generic gravitational systems could
not form singularities.

The event horizon changes the question.  Once a future-directed
timelike trajectory crosses it, no choice of aim or angular momentum
returns that trajectory to the exterior.  The singularity is not a
spatial target that can be missed; it is a future spacelike boundary.
Newtonian point-aiming arguments therefore do not prevent black-hole
formation.

\begin{classroomqa}
\audiencequestion{Is the same angular momentum that stabilizes a
cluster enough to prevent relativistic collapse?
\lecturetimestamp{01:52:29--01:56:31}}
\lecturerresponse{Angular momentum makes exact Newtonian point
collapse nongeneric and helps long-lived systems such as clusters
avoid immediate collapse.  But if a configuration lies inside its
event horizon, every future-directed causal path remains inside and
heads toward the singular future.  The old Newtonian argument treated
the singularity as a point one could miss and did not account for the
horizon.}
\end{classroomqa}

Pressure can resist collapse before a horizon forms.  Classical
Hamiltonian evolution preserves phase-space volume by Liouville's
theorem.  Quantum statistics and the uncertainty principle introduce
additional scales and, for fermions, degeneracy pressure.  These are
real sources of support, but they do not imply that an arbitrarily
massive object must remain outside its Schwarzschild radius.

\begin{classroomqa}
\audiencequestion{Is phase-space incompressibility a consequence of
Planck's constant, and can it always stop collapse?
\lecturetimestamp{01:56:33--01:58:22}}
\lecturerresponse{Liouville incompressibility is already a classical
theorem; it is not caused by \(\hbar\).  Quantum mechanics makes a
phase-space cell physically significant and can generate pressure,
but pressure is a local material response.  A sufficiently large
mass can have a Schwarzschild radius at a very low mean density, long
before microscopic compression becomes extreme.}
\end{classroomqa}

\section{Mean Density and Horizon Scale}

For a body of radius \(R\) and mean density \(\bar\rho\),
\begin{equation}
  M=\frac{4\pi}{3}\bar\rho R^3,
  \qquad
  R=\left(\frac{3M}{4\pi\bar\rho}\right)^{1/3}.
  \label{eq:l12-density-radius}
\end{equation}
At fixed \(\bar\rho\), the material radius grows as \(M^{1/3}\), while
the Schwarzschild radius grows as \(M\).  No matter how small the
density is, a sufficiently large mass satisfies
\begin{equation}
  R<2GM.
  \label{eq:l12-horizon-condition}
\end{equation}
At the threshold \(R=2GM\), the corresponding mean density is
\begin{equation}
  \bar\rho_{\mathrm{BH}}
  =\frac{3M}{4\pi(2GM)^3}
  =\frac{3}{32\pi G^3M^2}.
  \label{eq:l12-critical-density}
\end{equation}
The larger the black hole, the smaller its horizon-scale mean density.

\begin{figure}[tbp]
  \centering
  \includegraphics[width=0.84\textwidth]{lecture_12_figure_15.png}
  \caption{For fixed mean density,
  \(R\sim(M/\bar\rho)^{1/3}\), whereas
  \(r_{\mathrm s}=2GM\) grows linearly with mass.  A sufficiently
  large region therefore lies within its Schwarzschild radius even
  when the material is rarefied.}
  \lectureframe{02:02:10}
\end{figure}

This scaling explains why black-hole formation does not require
matter to begin at a universal enormous density.  For stellar masses,
material support and high density are central to the collapse
history.  For sufficiently large \(M\), the horizon condition itself
can occur at a modest average density.

\subsection{The cosmological comparison}

For approximately uniform vacuum energy density, the natural horizon
scale is cosmological.  In a spatially flat de Sitter approximation,
\begin{equation}
  H^2=\frac{8\pi G}{3}\rho_{\mathrm{vac}},
  \qquad
  R_{\mathrm{dS}}=H^{-1}.
  \label{eq:l12-de-sitter-scale}
\end{equation}
The Schwarzschild radius associated with the vacuum energy enclosed
inside a sphere becomes comparable to the sphere's radius precisely
at this scale.  The analogy is useful but must not be mistaken for
identity: a de Sitter observer is surrounded by a cosmological
horizon, not trapped inside a Schwarzschild black hole.

\begin{classroomqa}
\audiencequestion{Does the vacuum energy of our universe give a
Schwarzschild radius, and does that mean we live inside a black hole?
\lecturetimestamp{02:02:26--02:04:24}}
\lecturerresponse{The associated scale is the cosmological horizon,
of order the inverse Hubble parameter for vacuum-dominated expansion.
It is an inside-out analogue: distant objects recede and pass beyond
our causal horizon, whereas inside a black hole all future directions
lead inward toward a singularity.  A cosmological horizon is not the
interior of a Schwarzschild black hole.}
\end{classroomqa}

\section{What Classical Theory Cannot Continue}

The phrase end of time is shorthand, not a claim that the traveler's
clock takes infinitely long to arrive.  The proper time to \(r=0\) is
finite.  What ends is the classical spacetime predicted by the
Schwarzschild solution.  As the curvature, density, and tidal
gradients grow without bound, any fixed body of known matter is driven
beyond the domain where its familiar description applies.

\begin{classroomqa}
\audiencequestion{If the singularity is the end of time, how can an
observer actually reach it?
\lecturetimestamp{02:04:25--02:06:40}}
\lecturerresponse{The observer's proper time to the singularity is
finite.  Calling it an end means that the classical solution has no
regular future continuation there.  It is not meaningful to imagine
the observer stopping at fixed \(r=0\); fixed \(r\) was already
spacelike inside.  Curvature and material scales run beyond every
regime covered by the classical theory.}
\end{classroomqa}

There is also an operational way to see the loss of ordinary
structure.  Two neighboring constituents can influence one another
only if a light signal has time to pass between them.  As the
singularity approaches, the remaining proper time becomes shorter
than the required communication time over successively smaller
separations.  First separated atoms lose causal contact, then
electrons and nuclei, and eventually nuclear constituents.  The
classical picture is not a final microscopic theory; it says that no
finite bound structure can preserve its internal communication all
the way to the singular boundary.

\begin{figure}[tbp]
  \centering
  \includegraphics[width=0.84\textwidth]{lecture_12_figure_16.png}
  \caption{Near the spacelike singularity, neighboring worldlines
  have too little remaining proper time to exchange signals.  Bound
  structures lose causal communication on successively smaller
  scales.}
  \lectureframe{02:08:45}
\end{figure}

\begin{classroomqa}
\audiencequestion{Is the singularity merely irrelevant because no
signal from it can return to the exterior?
\lecturetimestamp{02:06:40--02:08:53}}
\lecturerresponse{It is causally irrelevant to an observer who remains
outside, but an infalling observer cannot turn it into an experiment
with a reported result.  Near the singularity, separated parts of the
apparatus lose causal contact before they could establish or transmit
a shared observation.}
\end{classroomqa}

The Rindler approximation cannot be used to infer this behavior.  It
was obtained under \(|r-r_{\mathrm s}|\ll r_{\mathrm s}\) and is valid
only near the regular horizon.  The global interior diagram and the
curvature divergence come from the full Schwarzschild geometry.

\begin{classroomqa}
\audiencequestion{Does the flat Rindler approximation remain valid
when we approach \(r=0\)?
\lecturetimestamp{02:08:53--02:09:43}}
\lecturerresponse{No.  The approximation was made by expanding about
\(r=r_{\mathrm s}\).  Moving toward \(r=0\) leaves that regime.  The
regularity of the horizon is established locally; the spacelike
singularity requires the full solution and has genuine curvature.}
\end{classroomqa}

\subsection{There Is No Classical Other Side of the Singularity}

At the horizon, a bad coordinate can be replaced by a regular one and
the metric extends smoothly.  At \(r=0\), the invariant curvature
diverges.  Substituting a negative number for the areal radius in the
Schwarzschild formula is therefore not an analytic continuation
through a regular surface.  Classical general relativity supplies no
spacetime beyond that singular boundary.

\begin{classroomqa}
\audiencequestion{Can a mathematician continue the metric to negative
\(r\) and call that the other side of the singularity?
\lecturetimestamp{02:09:48--02:12:34}}
\lecturerresponse{One may manipulate a formula beyond its physical
domain, but the curvature singularity blocks a regular spacetime
extension.  Unlike the horizon, there is no coordinate change that
removes the divergence and no observer can cross, return, and test a
putative continuation.  Questions about information and quantum
resolution require physics beyond this classical lecture.}
\end{classroomqa}

\section{Operational Tests at the Horizon}

Carry a meter stick across the horizon, with Alice at one end and
Shirley at the other.  Provided the stick is small compared with the
curvature scale, it remains an ordinary local object.  It does not
turn into a clock, and Alice's clock does not turn into a ruler.  Only
the Schwarzschild coordinate directions exchange causal roles.

\begin{figure}[tbp]
  \centering
  \includegraphics[width=0.84\textwidth]{lecture_12_figure_17.png}
  \caption{Alice and Shirley carry a finite meter stick through the
  horizon.  Their local apparatus remains regular; the dramatic
  exchange belongs to the Schwarzschild coordinate labels.}
  \lectureframe{02:13:00}
\end{figure}

\begin{classroomqa}
\audiencequestion{Does Alice's meter stick become a clock when it
crosses the horizon?
\lecturetimestamp{02:12:40--02:14:24}}
\lecturerresponse{Absolutely not.  The stick and its endpoints pass
through normally in a freely falling frame.  The symbol \(r\) no
longer labels a spacelike measuring direction inside, but naming a
coordinate \(r\) does not require it to be measured by rulers in every
region.  Local rods and clocks retain their ordinary roles.}
\end{classroomqa}

A second experiment uses reflected light.  Before crossing, Alice
sends a pulse outward toward Bob's reflector and receives the return.
As the emission point approaches the horizon, both legs of the trip
take longer in the exterior/Rindler description.  Motion of the
reflector relative to Alice redshifts the return.  Alice observes
progressively delayed and redshifted echoes, not a local image or
flash that announces the horizon at the instant of crossing.

\begin{figure}[tbp]
  \centering
  \includegraphics[width=0.84\textwidth]{lecture_12_figure_18.png}
  \caption{Alice sends pulses to Bob's reflector.  Successive return
  signals are delayed and redshifted as the horizon is approached,
  but no local discontinuity appears at crossing.}
  \lectureframe{02:16:20}
\end{figure}

\begin{classroomqa}
\audiencequestion{Can Alice locate the horizon by sending a signal to
Bob's reflector and watching the echo?
\lecturetimestamp{02:14:26--02:16:52}}
\lecturerresponse{She sees the round-trip time grow and the reflected
photons become strongly redshifted because Bob is accelerating away
in the local Rindler picture.  These changes are continuous.  No
return signal received at the crossing event identifies a local
horizon, and after crossing no future-directed signal reaches a
reflector that remains outside.}
\end{classroomqa}

\section{The Exterior Thin-Shell Description}

From the viewpoint of a Schwarzschild observer, infalling layers
approach the horizon asymptotically and become ever thinner and more
redshifted.  It is therefore permissible within that exterior
description to picture the accumulated mass as sedimentary layers
concentrated just outside the horizon.  This is a description tied to
the exterior time slicing, not a material membrane encountered by a
freely falling observer.

\begin{figure}[tbp]
  \centering
  \includegraphics[width=0.84\textwidth]{lecture_12_figure_19.png}
  \caption{The exterior slicing represents infalling material as
  increasingly thin layers near the horizon.  The surrounding field
  records the enclosed mass even though no signal escapes from the
  interior.}
  \lectureframe{02:18:45}
\end{figure}

\begin{classroomqa}
\audiencequestion{If gravitons cannot escape the black hole, how can
the exterior gravitational field persist?
\lecturetimestamp{02:17:13--02:18:58}}
\lecturerresponse{The picture of force as particles shuttling back and
forth should not be taken literally.  Classically, the exterior
geometry is the solution of the field equations with the conserved
mass as boundary data.  If one insists on an exchange-particle
heuristic, the exterior slicing places the effective source in the
ever-thinner layers just above the horizon; no graviton must carry
news outward from the singular interior.}
\end{classroomqa}

\section{Rotating Black Holes}

Astrophysical black holes generally rotate.  Their stationary
exterior is described by the Kerr solution rather than the
Schwarzschild solution.  In units \(c=1\), a Kerr horizon exists only
when
\begin{equation}
  |J|\leq GM^2,
  \qquad
  a_*=\frac{J}{GM^2},
  \qquad
  |a_*|\leq1.
  \label{eq:l12-kerr-bound}
\end{equation}
Rotation flattens the horizon and introduces frame dragging, but the
dimensionless spin is bounded.  A black hole therefore cannot become
an arbitrarily flattened material disk.

\begin{classroomqa}
\audiencequestion{Could a rotating black hole contain a stable binary
system orbiting inside it, and would gravitons from that system have
to escape?
\lecturetimestamp{02:19:00--02:20:22}}
\lecturerresponse{No stationary internal binary can remain hidden
inside the horizon and keep orbiting forever.  Once inside, its
future-directed motion belongs to the black-hole interior and ends at
the singular region.  The exterior Kerr field is characterized by
global charges such as \(M\) and \(J\); it does not require signals or
gravitons to escape from an enduring mechanical system inside.}
\end{classroomqa}

\section{Summary}

The Schwarzschild horizon is locally simple.  Proper distance from
the horizon and the rescaled time \(t/(4GM)\) turn its metric into the
Rindler metric of flat spacetime in accelerated coordinates.  This
explains both sides of the familiar account: a freely falling
observer crosses in finite proper time without encountering local
singular curvature, while a supported exterior observer receives
signals that are delayed and redshifted without bound.

The interior is not merely another Rindler wedge.  Fixed \(r\) becomes
spacelike, the singularity at \(r=0\) is a spacelike future boundary,
and the curvature invariant grows as \(r^{-6}\).  Angular momentum,
pressure, and phase-space constraints can delay or prevent a system
from entering its Schwarzschild radius, but after a horizon forms they
do not provide an outward future.  The vacuum field equations
determine the exterior from symmetry and conserved mass, the
horizon-scale mean density decreases as \(M^{-2}\), and the classical
description finally ends where causal communication and curvature
scales run beyond its domain.  Rotating black holes modify the
geometry but preserve the central lesson: a horizon is regular
locally and decisive globally.
